\section{Índice de Efetividade da Gestão Municipal}
Concebido originalmente em 2014 pelo Tribunal de Contas do Estado de São Paulo (TCE-SP), o Índice de 
Efetividade da Gestão Municipal (IEGM) foi posteriormente incorporado e aperfeiçoado pelo Instituto 
Rui Barbosa (IRB) - entidade voltada ao desenvolvimento de estudos e pesquisas para o sistema de 
controle externo \cite{iegm_ceara_2022}. Com o intuito de uniformizar a avaliação das administrações 
públicas locais, o índice foi selecionado em 2016 pela Rede Nacional de Indicadores Públicos 
(Indicon) para aplicação em âmbito nacional \cite{iegm_ceara_2022}.

O indicador foi projetado para mensurar o nível de eficácia e desenvolvimento das 
administrações municipais brasileiras. O IEGM engloba sete áreas fundamentais: educação, saúde, 
gestão fiscal, planejamento, meio ambiente, defesa civil e governança de TI \cite{lima_influencia_2023}.

O resultado geral do IEGM varia de 0 a 100, no qual a pontuação máxima representa a aderência 
plena do município às normas e aos controles estabelecidos. Para o cálculo do índice consolidado, 
aplica-se uma métrica de ponderação que reflete a relevância de cada dimensão específica. Assim, a 
composição do IEGM é estruturada pelas seguintes proporções:
\cite{lima_influencia_2023}.


\begin{itemize}
  \item Educação (i-Educ, 20\%)
  \item Saúde (i-Saúde, 20\%)
  \item Gestão Fiscal (i-Fiscal, 20\%)
  \item \revisar{Planejamento (i-Planejamento, 20\%)}
  \item Meio Ambiente (i-Amb, 10\%)
  \item Proteção dos Cidadãos/Defesa Civil (i-Cidade, 5\%)
  \item Governança de TI (i-Gov TI, 5\%)
\end{itemize}

O Índice conta com a adesão de todos os 27 Tribunais de Contas dos estados e do Distrito Federal, 
consolidando-se como um instrumento nacional de avaliação das administrações públicas locais \cite{araujo_gestoes_municipais_2021}. 
A operacionalização do índice ocorre por meio do Termo de Adesão à Rede de Indicadores Nacionais, 
firmado entre o Instituto Rui Barbosa (IRB) e os órgãos de contas estaduais e distrital. A partir 
desse acordo, as informações prestadas pelos municípios são coletadas e posteriormente validadas 
in loco ou auditadas pelos respectivos tribunais, garantindo a consistência dos dados do índice \cite{araujo_gestoes_municipais_2021}.

\subsection{Índice de Governança de TI (I-Gov TI)}

\revisar{No contexto da transformação digital do setor público, a governança de Tecnologia da Informação 
posiciona-se como elemento indispensável para a otimização dos serviços estatais. Ao impulsionar a 
transparência, a integridade e a eficiência das entregas governamentais, o uso estratégico da TI 
reflete diretamente na efetividade da gestão pública. O I-Gov TI surge, portanto, como um indicador 
estruturado para diagnosticar e induzir o aprimoramento da governança tecnológica municipal, em 
estrita observância aos princípios da moralidade e da eficiência administrativa \cite{silva2025implementaccao}.}

\revisar{Como componente estruturante do IEGM, o i-Gov TI mensura o alinhamento entre a gestão de TI e as 
demandas coletivas. Sua avaliação incide sobre os pilares da governança digital municipal: 
diretrizes para uso de TIC, protocolos de segurança da informação, desenvolvimento de competências 
da equipe interna e fomento à transparência governamental \cite{construcao_iegm_2018}.}


\revisar{A aferição desses aspectos reflete de forma direta a capacidade dos municípios em atender às 
exigências normativas e em incorporar as melhores práticas de gestão pública. Cumpre destacar que 
toda essa estrutura avaliativa fundamenta-se rigorosamente nos diplomas legislativos que orientam, 
normatizam e direcionam a boa governança de Tecnologia da Informação \cite{CEPCP_2025}.}

\revisar{Nessa perspectiva, o i-Gov TI alinha-se a tais exigências ao verificar a participação formal da 
equipe de Tecnologia da Informação nos processos licitatórios do setor. A exigência do Plano Diretor
de Tecnologia da Informação (PDTI) para a elaboração de um Estudo Técnico Preliminar (ETP) 
consistente posiciona o planejamento como instrumento central de economicidade e conformidade 
contratual. Corroborando essa diretriz, a jurisprudência do Tribunal de Contas da União (TCU) aponta 
recorrentemente a deficiência no planejamento como causa primária de contratações ineficazes. A 
ausência de planejamento estratégico em TI, portanto, compromete a eficiência dos investimentos 
públicos e expõe os municípios a riscos normativos e à responsabilização de seus gestores \cite{CEPCP_2025}.}

\revisar{Além do respaldo normativo e jurisprudencial, a avaliação do i-Gov TI apoia-se em frameworks 
internacionais de governança de TI. O COBIT, por exemplo, estrutura práticas e processos voltados ao 
alinhamento estratégico da tecnologia com os objetivos institucionais e a geração de valor público. 
A dimensão i-Gov TI incorpora essa perspectiva ao averiguar se os municípios possuem planejamento 
estratégico de TI integrado ao planejamento governamental global, instigando a superação de 
iniciativas isoladas \cite{CEPCP_2025}.}